Les enjeux que présente l'accessibilité aux collections des institutions patrimoniales ont donc trait à notre capacité à avoir accès à l'ensemble des collections, sans avoir connaissance des pratiques documentaires utilisées. En l'état actuel, il est difficile d'avoir une utilisation autonome des collections sans avoir de connaissances sur l'histoire du traitement des collections. Nous avons identifié les raisons de ce manque de visibilité sur l'ensemble des collections comme étant à la fois un mélange de plusieurs politiques documentaires facilement identifiables comme des périodes clés dans l'histoire de l'institution, une accumulation de changements mineurs non documentés et une adaptation des pratiques de traitement documentaire à l'informatisation des pratiques documentaires. Il en résulte une accessibilité qui repose principalement sur la transmission orale des connaissances approfondies des documentalistes. Afin d'offrir une meilleure accessibilité de ces documents, nous avons proposé la conception d'un outil d'appui aux outils de consultation des bases de données, proposant une cartographie du traitement documentaire des collections afin de fournir à l'utilisateur une représentation graphique compréhensible des données qu'il exploite, ainsi qu'une plateforme à même d'accueillir et de conserver les connaissances des documentalistes.

Cela dit, nous avons également constaté, tout au long de cette analyse, que certains outils sont en cours de développement. Ces outils offrent des solutions aux défauts des bases de données. Le lac de données, notamment, a effectué un travail important sur la qualité des données lors de la migration des bases de deux services patrimoniaux de l'INA. De plus, si nous avons constaté l'impact de l'enchaînement de différentes politiques documentaires, force est de constater que l'utilisation de l'intelligence artificielle est en passe de devenir la prochaine évolution des politiques documentaires. Seul l'avenir ne peut nous dire si cette nouvelle évolution sera avantageuse pour l'accessibilité aux collections ou si, au contraire, le bruit que provoquera ces données générées n'accentuent la confusion. 